\section{A finite obstruction to collapsing the fourth implicational parameter}

Write the unrestricted Lukasiewicz form as
\[
 R(p,q,r,u)=((p\to u)\to q)\to((q\to p)\to(r\to p)).
\]
The PA emitter retains the restricted family $u\le q$ on numerical codes.
A tempting simplification is to set $u=q$ throughout, eliminating the
independent fourth parameter. This family alone, even with modus ponens
and schematic substitution, is incomplete.

To prove this, take the set $\{0,1,2\}$, designate only $2$, and interpret
implication by the following matrix (antecedents index rows):
\[
\begin{array}{c|ccc}
\to&0&1&2\\\hline
0&1&1&2\\
1&1&2&2\\
2&1&1&2
\end{array}
\]
Every instance of
\[
 ((p\to q)\to q)\to((q\to p)\to(r\to p))
\]
has value $2$, as direct evaluation of the $27$ triples shows.
Modus ponens preserves designation: when its first premise has value $2$,
row $2$ has designated output only in column $2$. But $p\to p$ has
value $1$ at $p=0$. Hence identity is not derivable from this specialized
basis.

The supplied finite checker constructs this matrix by exhaustive search
and then evaluates all its cases. It is a mathematical obstruction to the
specialized implicational basis, not a claim that all the additional
quantifier and set-theoretic axioms of the complete numerical calculus
cannot restore the missing logic. Such a restoration would require a new
proof before using the specialization in a consistency machine.
